\documentclass[10pt,a4paper]{amsart}
\usepackage[margin=19mm]{geometry}
\usepackage{amsmath,amssymb,mathtools}
\usepackage[scr=rsfs,cal=euler]{mathalfa}
\usepackage{xfrac}
\usepackage[unicode,hidelinks,bookmarks=false]{hyperref}
\newcommand{\ie}{i.e.\ }
\newcommand{\cf}{cf.\ }
\newcommand{\resp}{respectively\ }
\newcommand{\Subsection}{\subsection}
\providecommand{\nobreakdash}{\nobreak\hbox{-}}
\setlength{\emergencystretch}{2em}
\multlinegap0pt
\usepackage{mathtools}
\usepackage{amssymb}
\usepackage{bm}
\usepackage{enumerate}
\usepackage[shortlabels]{enumitem}
\usepackage{tikz}
\newtheorem{prop}{Proposition}
\newcommand{\End}{\operatorname{End}}
\newcommand{\Env}{\operatorname{Env}}
\newcommand{\Hom}{\operatorname{Hom}}
\newcommand{\PaRB}{\mathsf{PaRB}}
\newcommand{\bT}{\mathbf{T}}
\newcommand{\coev}{\operatorname{coev}}
\newcommand{\ev}{\operatorname{ev}}
\newcommand{\id}{\operatorname{id}}
\newcommand{\ob}{\operatorname{ob}}
\newcommand{\twist}{\tau}
\title{Grothendieck-Teichm\"uller Symmetries of Cyclic Operads and Tangles}
\author{Marcy Robertson, Chandan Singh}
\date{Source: arXiv:2511.05911v3 (CC BY 4.0)}
\begin{document}
\maketitle

\noindent\emph{Source and licence.} Grothendieck-Teichm\"uller Symmetries of Cyclic Operads and Tangles. Version arXiv:2511.05911v3 (CC BY 4.0), distributed by arXiv under the Creative Commons Attribution 4.0 International licence, http://creativecommons.org/licenses/by/4.0/, as recorded in the arXiv licence metadata for this exact version and shown on the abstract page https://arxiv.org/abs/2511.05911v3. Full licence notice: \texttt{licenses/arxiv-2511.05911-CC-BY-4.0.txt}.\par\smallskip

\noindent\textit{Editorial scope.} Counted unit: the article's prop endomorphisms (source0.tex lines 6907--6911). The article's complete original proof of it is delivered with it (source0.tex lines 6913--7031, one \texttt{proof} environment of 119 lines). No mathematics is written, repaired or re-derived: the statement and the proof are the article's own lines, in the article's own notation, and no retained line is substituted or re-flowed.  Retained uncounted as the source-local context the proof needs: lines 852--861; lines 3497--3505; lines 3508--3514.  One declared adaptation: exactly 1 ancillary strings inside the retained spans are omitted (image-inclusion commands and, where the source repeats a label, the repeated label definitions) and no other line differs from the source; the figure environments, with their placement, captions and labels, are kept, and the package records the omitted strings in fidelity receipts bound to the affected bodies.  No retained line contains a citation, so the wrapper carries no bibliography and no bibliography entry.  The retained lines contain the article's own diagram code, which the wrapper typesets with the standard tikz packages; no image file is delivered and the package declares no asset dependency.  Editorial formatting only, in addition: an AMS \texttt{amsart} preamble that restates the article's own declarations and macros used by the retained lines, namely the environments \texttt{prop} on separate counters, together with the standard packages those lines need.  The retained environments print in this document as Proposition 1 in order of appearance, whereas the article prints the same items with the numbering of its own full document.  The complete native source of the article as distributed in the arXiv e-print for this exact version is bundled unchanged as \texttt{sources/188689/source0.tex}.
\begin{equation}
\label{eq: transpose of a morphism}
f^*
\coloneqq
(\ev_Y\otimes\id_{X^*})
\cdot
(\id_{Y^*}\otimes f\otimes\id_{X^*})
\cdot
(\id_{Y^*}\otimes\coev_X).
\end{equation}

\begin{equation}
\label{eq: cyclic action PaRB generators}
s_0^*\twist = \twist,
\qquad
s_0^*\beta_{12}
= \bigl(1_\mu\circ_2\twist^{-1}\bigr)\cdot\beta_{12}^{-1},
\qquad
s_0^*\alpha_{123} = \alpha_{231}^{-1}.
\end{equation}

\begin{figure}[h]
    \centering
    
    \caption{The cyclic action on the generating morphisms of $\PaRB$,
as given in \eqref{eq: cyclic action PaRB generators}.}
    \label{fig:cyclic_action_on_PaRB}
\end{figure}

\begin{prop}
\label{prop: metric prop action through cyclic endomorphisms}
The $\Env(\PaRB)$-action on $q\bT'$ extends to an action of
$\Pi(\PaRB^{\mathrm{cyc}})$ through $\End_{q\bT'}$.
\end{prop}

\begin{proof}
We first verify that the functors and natural transformations
described above define the $\Env(\PaRB)$-action.
For $\omega\in\ob(\PaRB(n))$ and
$\omega'\in\ob(\PaRB(m))$, an element of the coend describing
substitution is represented by morphisms
\[
u\colon w_0\longrightarrow\omega(\ldots,w^*,\ldots),
\qquad
v\colon w^*\longrightarrow\omega'(\ldots).
\]
Send this representative to $\omega(\ldots,v,\ldots)\cdot u$.
The coend relation identifies the two ways of applying a morphism
of the intermediate word. Conversely, a morphism
$w_0\to\omega(\ldots,\omega'(\ldots),\ldots)$ is represented
by taking $w=\omega'(\ldots)^*$ and
$v=\id_{\omega'(\ldots)}$, using the double-dual identification.
The coend relations identify every representative with one of
this form, giving a bijection.
Thus composition agrees with operadic substitution.

The relations on generating morphisms hold by the given
$\PaRB$-action. Taking Cartesian products on separate lists of
variables and permuting those variables gives the corresponding
$\Env(\PaRB)$-action. Evaluating at dual input words and using
$w_i^{**}\cong w_i$ recovers the original $\PaRB$-action.

To extend this action to the metric prop, assign both $\ev$
and $\coev$ the pairing
$(w,w')\mapsto\Hom(w,(w')^*)$, with its variables regarded
as inputs for $\ev$ and outputs for $\coev$.
We verify the symmetry, snake, and cyclic-invariance relations.

For $h\colon w\to(w')^*$, its transpose, defined in
\eqref{eq: transpose of a morphism}, is
$h^*\colon(w')^{**}\to w^*$.
The symmetry $\Hom(w,(w')^*)\cong\Hom(w',w^*)$ sends $h$
to the composite
\[
w'\xrightarrow{\cong}(w')^{**}\xrightarrow{h^*}w^*.
\]
Applying this construction twice returns $h$ under the
double-dual identifications.

For the snake relations, fix words $w_0,w_1$. An element of
the first snake composite evaluated at $(w_0;w_1)$ is
represented by a chain
\[
w_0\xrightarrow{h_1}w_2^*
\xrightarrow{h_2}w_3
\xrightarrow{h_3}w_4^*
\xrightarrow{h_4}w_1^*.
\]
Composition sends its class to
$h_4\cdot h_3\cdot h_2\cdot h_1\in\Hom(w_0,w_1^*)$.
An inverse represents $u\colon w_0\to w_1^*$ by taking
$w_2=w_4=w_0^*$, $w_3=w_0$, the first three morphisms to
be identities, and the last to be $u$.
The coend relations identify every chain with this representative.
The second snake composite gives the same bijection after
transposing the factors. Hence both composites are isomorphic
to the identity operation
$(w_0;w_1)\mapsto\Hom(w_0,w_1^*)$
in $\End_{q\bT'}(1,1)$.

For cyclic invariance, let $H$ be a contravariant functor of
$n+1$ variables, with $n\geq1$. Contracting the identity
factors in the image of the cyclic-invariance relation gives
\[
\begin{aligned}
&\int^{w,w'}
\Hom(w_0,w)\times
H(w',w,w_1,\ldots,w_{n-1})\times
\Hom(w_n,w')\\
&\qquad\cong
H(w_n,w_0,\ldots,w_{n-1})
=(z_{n+1}^*H)(w_0,\ldots,w_n).
\end{aligned}
\]
A representative consists of $a\colon w_0\to w$,
$b\colon w_n\to w'$, and
$h\in H(w',w,w_1,\ldots,w_{n-1})$.
The displayed map sends its class to
$H(b,a,\id,\ldots,\id)(h)$.
Its inverse takes $w=w_0$, $w'=w_n$, and $a,b$ to be
identities. Thus the images of $\ev$ and $\coev$ implement
the required cyclic permutation.

For the functor assigned to $\omega$, duality gives
\[
\Hom\bigl(w_0,\omega(w_1^*,\ldots,w_n^*)\bigr)
\cong
\Hom\bigl(
\emptyset,w_0^*\otimes\omega(w_1^*,\ldots,w_n^*)
\bigr).
\]
Under this identification, changing the output is implemented
by bending strands through cups and caps. On objects, this
reroots the planar tree representing $\omega$. On generating
morphisms, the bending comparisons give
\[
s_0^*\twist=\twist,\qquad
s_0^*\beta_{12}
=(1_\mu\circ_2\twist^{-1})\cdot\beta_{12}^{-1},
\qquad
s_0^*\alpha_{123}=\alpha_{231}^{-1},
\]
as in \eqref{eq: cyclic action PaRB generators} and
Figure~\ref{fig:cyclic_action_on_PaRB}.
These comparisons respect gluing and relabelling by cup--cap
cancellation and ribbon coherence. Since the displayed morphisms
generate $\PaRB$, they establish compatibility with its cyclic
structure.

The assignments therefore preserve the defining relations of
$\Pi(\PaRB^{\mathrm{cyc}})$. The coend and duality comparisons
are compatible with iterated composition and tensor product,
giving the asserted extension.
\end{proof}

\end{document}
